\documentclass[12pt, leqno, letterpaper]{article}


\renewcommand{\baselinestretch}{1.0}
\usepackage{makecell}
\usepackage{booktabs}

\usepackage{amsmath}
\usepackage{amssymb}
\usepackage{enumitem}
\usepackage{graphicx}
\usepackage{comment}
\usepackage{scrextend}
\usepackage{pgffor}
\usepackage{ifxetex}


\ifxetex
  \usepackage[T1]{fontenc}
  \usepackage[no-math]{fontspec}

  \setmainfont[Scale=1.0]{Kepler Std Light}

  \usepackage[mtpcal,mtphrd,mtpfrak,slantedGreek,noamssymbols,
  	      subscriptcorrection,zswash]{mtpro2}

  \DeclareMathSizes{12}{12}{9}{7}
\else
  \usepackage[T1]{fontenc}
  \usepackage{lmodern}
  % to match mt2pro
  \newcommand{\mbf}{\mathbf} 

\fi
 

\usepackage{bm}

\usepackage[scr=boondox, bb = boondox,scrscaled=1]{mathalfa}

\newcommand{\diag}{\textbf{diag}}   
\newcommand{\var}{\textbf{var}}   
\newcommand{\Exp}{\textbf{E}}   
\newcommand{\tsum}{{\textstyle\sum\nolimits}}

\newcommand{\ep}{\epsilon}   
\newcommand{\tq}[1]{{\textquotedblleft #1\textquotedblright}}
\newcommand{\s}[1][1]{\hspace{#1pt}}

\newcommand{\np}{n/\s[-1.5] p}
\newcommand{\qeig}{\calS}
\newcommand{\seig}{\scrs}
\newcommand{\Id}{I}
\newcommand{\mvw}{\hw_\textsc{mv}}
\newcommand{\sfw}{\hw_\textsc{cp}}
\newcommand{\hsig}{\hat{\sigma}}
\newcommand{\hSig}{\bm{\hat}{\Sigma}}
\newcommand{\bSig}{{\Sigma}}
\newcommand{\ip}[2]{\langle #1, #2 \rangle}

\newcommand{\hw}{\hat{w}}
\newcommand{\req}[1]{(\ref{#1})}
\newcommand{\tru}{\scrV}
\newcommand{\nes}{\mathbb{1}}

\foreach \x in {A,B,...,Z,a,b,...,z} 
{
  \expandafter\xdef\csname cal\x\endcsname{\noexpand 
	\ensuremath{\noexpand\mathcal{\x}}}
  \expandafter\xdef\csname scr\x\endcsname{\noexpand 
	\ensuremath{\noexpand\mathscr{\x}}}
  \expandafter\xdef\csname bb\x\endcsname{\noexpand 
	\ensuremath{\noexpand\mathbb{\x}}}
  \expandafter\xdef\csname rm\x\endcsname{\noexpand 
	\ensuremath{\noexpand\mathrm{\x}}}
  \expandafter\xdef\csname bf\x\endcsname{\noexpand 
	\ensuremath{\noexpand\mbf{\x}}}
}

\begin{document}


\section{Model}

For a factor return $f \in \bbR^q$ and specific return
$\ep \in \bbR^p$, the excess security return,
\begin{align}
 y =  B f + \ep \s .
\end{align}
with $f$ and $\ep$ uncorrelated. 
$\var(y) = B\var(f) B^\top + \var(\ep)$ 
with $\var(\ep)$ diagonal.
\begin{align} 
 \var(f) = \left( \begin{array}{ccccccc} 
$256$ & $17$ & $-9$ & $51$ & $19$ & $6$ & $1$ \\ 
$17$ & $16$ & $0$ & $0$ & $0$ & $0$ & $0$ \\ 
$-9$ & $0$ & $4$ & $0$ & $0$ & $0$ & $0$ \\ 
$51$ & $0$ & $0$ & $400$ & $0$ & $0$ & $0$ \\ 
$19$ & $0$ & $0$ & $0$ & $225$ & $0$ & $0$ \\ 
$6$ & $0$ & $0$ & $0$ & $0$ & $100$ & $0$ \\ 
$1$ & $0$ & $0$ & $0$ & $0$ & $0$ & $25$ \\ 
\end{array} \right) 
 \end{align} 
\begin{align} 
 \text{corr}(f) = \left( \begin{array}{ccccccc} 
$1.00$ & $0.28$ & $-0.30$ & $0.16$ & $0.08$ & $0.04$ & $0.02$ \\ 
$0.28$ & $1.00$ & $-0.11$ & $0.00$ & $0.00$ & $0.00$ & $0.00$ \\ 
$-0.30$ & $-0.11$ & $1.00$ & $0.00$ & $0.00$ & $0.00$ & $0.00$ \\ 
$0.16$ & $0.00$ & $0.00$ & $1.00$ & $0.00$ & $0.00$ & $0.00$ \\ 
$0.08$ & $0.00$ & $0.00$ & $0.00$ & $1.00$ & $0.00$ & $0.00$ \\ 
$0.04$ & $0.00$ & $0.00$ & $0.00$ & $0.00$ & $1.00$ & $0.00$ \\ 
$0.02$ & $0.00$ & $0.00$ & $0.00$ & $0.00$ & $0.00$ & $1.00$ \\ 
\end{array} \right) 
 \end{align}Factor volatilities $\sigma_f = \sqrt{\diag(\var(f))}$ are given by, 
 \begin{align} 
 \sigma_f = \left( \begin{array}{ccccccc} 
$16.0$ & $4.0$ & $2.0$ & $20.0$ & $15.0$ & $10.0$ & $5.0$ \\ 
\end{array} \right) 
 \end{align}with expectations, \begin{align} 
 \Exp(f) = \left( \begin{array}{ccccccc} 
$6.00$ & $3.48$ & $4.44$ & $0.00$ & $0.00$ & $0.00$ & $0.00$ \\ 
\end{array} \right) 
 \end{align}Population means $\alpha = \Exp(y) = B \Exp(f) + \Exp(\ep)$ have the averages,\begin{align*} 
p^{-1} \tsum_{i=1}^p \alpha_i = 6.23\quad  p^{-1} \tsum_{i=1}^p (B\Exp(f))_i = 5.98\quad p^{-1} \tsum_{i=1}^p \Exp(\ep_i) = 0.25
 \end{align*}Population volatilities $\sigma_X =\sqrt{\diag(\var(X))}$ are given by, 
 \begin{align*} 
p^{-1} \tsum_{i=1}^p (\sigma_y)_i = 41.27\quad  p^{-1} \tsum_{i=1}^p (\sigma_{Bf})_i = 21.05\quad p^{-1} \tsum_{i=1}^p (\sigma_{\ep})_i = 34.95 \, .
 \end{align*}

\begin{figure}[htp!]
\centering
\includegraphics[width=0.32\textwidth]{img/vols_alph.pdf}
\includegraphics[width=0.32\textwidth]{img/Bvol_Bret.pdf}
\includegraphics[width=0.32\textwidth]{img/svol_sret.pdf}
\caption{Scatter plots of $\alpha, B \s \Exp(f), \Exp(\ep)$
versus the corresponding volatilities.
}
\label{fig:hist}
\end{figure}


\begin{figure}[htp!]
\centering
%\includegraphics[width=0.475\textwidth]{img/block_orig.pdf}
\includegraphics[width=0.475\textwidth]{img/load_hist.pdf}
\includegraphics[width=0.475\textwidth]{img/block_post.pdf}
\caption{Left: Histograms of dense factor exposures
in $\Xi$. Right: Reordered sparse factor component 
$\scrS$ of
$\Xi\Xi^\top = \text{\tq{dense}} + \scrS$}
\label{fig:exp}
\end{figure}

\begin{figure}[htp!]
\centering
\includegraphics[width=0.475\textwidth]{img/frontier500.pdf}
\includegraphics[width=0.475\textwidth]{img/frontier1000.pdf}
\includegraphics[width=0.475\textwidth]{img/frontier2000.pdf}
\includegraphics[width=0.475\textwidth]{img/frontier3000.pdf}
\caption{Effient (population) frontiers for each $p$.}
\label{fig:pf}
\end{figure}



\begin{comment}
\begin{figure}[htp!]
\centering
\includegraphics[width=0.49\textwidth]{img/Bret_hist.pdf}
\includegraphics[width=0.49\textwidth]{img/sret_hist.pdf}
\caption{\textup{Left panel:} Histogram of $B \Exp(f)$.
\textup{Right panel:} Histogram of  $\Exp(\ep)$.
%Histogram of asset specific volatilities
% (i.e., square roots of the diagonal entries of $\Sv$).
}
\label{fig:hist}
\end{figure}



\begin{figure}[htp!]
\centering
\includegraphics[width=0.45\textwidth]{img/svol_sret.pdf}
\includegraphics[width=0.45\textwidth]{img/svol_hist.pdf}
\caption{Histogram of specific vols
 (i.e., square roots of the diag. entries of $\Sv$).}
\label{fig:hist}
\end{figure}


\begin{figure}[htp!]
\centering
\includegraphics[width=0.45\textwidth]{img/Bvol_Bret.pdf}
\includegraphics[width=0.45\textwidth]{img/Bvol_hist.pdf}
\caption{Histogram of specific vols
 (i.e., square roots of the diag. entries of $\Sv$).}
\label{fig:hist}
\end{figure}

\end{comment}
\clearpage

\section{Results}

For $\hSig$ and estimate of the population covariance
$\bSig$, we solve
\begin{align} \label{mqp}
\begin{aligned}
 \min_{x \in \bbR^p}\s  \s \frac{1}{2} \s &\ip{x}{\hSig x} \\
 \ip{m}{w} &\ge \mu \s ,
\\ \ip{\rme}{w} &= 1 \s ,
\end{aligned}
\end{align}
where the constraint is nonbining in the (rare) cases when the
minimum variance portfolio return exceeds $\mu$.
The constraint $m$ is an estimate of $\alpha = \Exp(y)$.

Let $\hw$ be the solution of $\req{mqp}$.
The optimal portfilio $w$ solves $\req{mqp}$ with $(\bSig,\alpha)$ 
replacing $(\hSig,m)$. We work with the following metrics.
\begin{align}
\tru &= \sqrt{\ip{\hw}{\bSig \hw}} \s[32]
\big(\text{True Volatility (TV)}\big) \\
\hat{\sigma} &= \sqrt{\ip{\hw}{\hSig \hw}} \s[32]
\big(\text{Estimated Volatility (EV)}\big) \\
\sigma &= \sqrt{\ip{w}{\bSig w}} \s[32]
\big(\text{Optimal Volatility}\big) \\
\Pi &= \frac{\hat{\sigma}}{\tru} \s[64]
\big(\text{Precision}\big)  \\ 
\scrR &= \frac{\ip{\hw}{\alpha} - r_f}{\tru} \s[32]
\big(\text{True Sharpe Ratio (TSR)}\big) \\
\hat{\rho} &= \frac{\ip{\hw}{m} - r_f}{\hsig} \s[32]
\big(\text{Estimated Sharpe Ratio (ESR)}\big) 
\\\rho &= \frac{\ip{w}{\alpha} - r_f}{\sigma} \s[32]
\big(\text{Optimal Sharpe Ratio}\big) 
\end{align}
where $r_f$ is the risk free rate.


The portfilio $\hw$ decomposes into a minimum variance (MV)
portfolio $\mvw$ and the characteristic portfilio (CP)
portfolio $\sfw$,  i.e., 
\[ \hw = \mvw + \sfw  \]
This allows for a return to $\hw$ to be decomposed.
\begin{align}
&\ip{\hw}{m} \s[32] \big(\text{Estimated Return (ER)}\big) \\
&\ip{\hw}{\alpha} \s[32] \big(\text{True Return (TR)}\big) \\
&\ip{\mvw}{m} \s[32] \big(\text{Estimated MV Return (MVER)}\big) \\
&\ip{\sfw}{m} \s[32] \big(\text{Estimated CP Return (CPER)}\big) \\
&\ip{\mvw}{\alpha} \s[32] \big(\text{True MV Return (MVER)}\big) \\
&\ip{\sfw}{\alpha} \s[32] \big(\text{True CP Return (CPER)}\big) 
\end{align}


\begin{table}[hpt!] 
\begin{center} 
\begin{tabular}{rccccccc} 
\toprule 
$p$ & \textsc{opt}  & \textsc{pca tv} & \textsc{pca ev} & $\calP$ & \textsc{jsm tv} & \textsc{jsm ev} & $\calP$ \\ 
\midrule 
$500$ & $6.57$ & $9.1$ & $6.03$ & $0.66$ & $8.08$ & $8.35$ & $1.03$ \\ 
$1000$ & $4.92$ & $8.15$ & $4.48$ & $0.55$ & $6.35$ & $6.55$ & $1.03$ \\ 
$2000$ & $3.5$ & $7.26$ & $3.17$ & $0.44$ & $4.71$ & $4.71$ & $1.0$ \\ 
$3000$ & $2.87$ & $7.02$ & $2.61$ & $0.37$ & $4.0$ & $3.97$ & $0.99$ \\ 
\bottomrule 
\end{tabular} 
\end{center} 
\caption{Portfolio volatility statistics ($n = 125$, $\mu = 8.5$,  10 simulations).} 
\label{tab:vol} 
\end{table} 

 

\begin{table}[hpt!] 
\begin{center} 
\begin{tabular}{rccccccc} 
\toprule 
$p$ & \textsc{opt}  & \textsc{pca tsr} & \textsc{pca esr} & \calD & \textsc{jsm tsr} & \textsc{jsm esr} & \calD \\ 
\midrule 
$500$ & $1.29$ & $0.72$ & $1.7$ & $2.36$ & $0.85$ & $1.21$ & $1.43$ \\ 
$1000$ & $1.73$ & $0.81$ & $2.34$ & $2.88$ & $1.1$ & $1.58$ & $1.43$ \\ 
$2000$ & $2.43$ & $0.92$ & $3.32$ & $3.61$ & $1.52$ & $2.22$ & $1.46$ \\ 
$3000$ & $2.96$ & $0.96$ & $4.06$ & $4.25$ & $1.82$ & $2.66$ & $1.46$ \\ 
\bottomrule 
\end{tabular} 
\end{center} 
\caption{True/Estimated Sharp Ratio (SR) statistics ($n = 125$, $\mu = 8.5$,  10 simulations).} 
\label{tab:SR} 
\end{table} 

 

\begin{table}[hpt!] 
\begin{center} 
\begin{tabular}{rccccccc} 
\toprule 
$p$ & \textsc{opt}  & \textsc{pca tr} & \textsc{pca er} & \calD & \textsc{jsm tr} & \textsc{jsm er} & \calD \\ 
\midrule 
$500$ & $8.5$ & $6.54$ & $8.5$ & $1.3$ & $6.85$ & $8.5$ & $1.24$ \\ 
$1000$ & $8.5$ & $6.6$ & $8.5$ & $1.29$ & $7.01$ & $8.5$ & $1.22$ \\ 
$2000$ & $8.5$ & $6.68$ & $8.5$ & $1.27$ & $7.18$ & $8.5$ & $1.19$ \\ 
$3000$ & $8.5$ & $6.69$ & $8.5$ & $1.27$ & $7.28$ & $8.5$ & $1.17$ \\ 
\bottomrule 
\end{tabular} 
\end{center} 
\caption{True/estimated return statistics -- True Return (TR) and Estimated Return (ER) -- ($n = 125$, $\mu = 8.5$,  10 simulations).} 
\label{tab:RS} 
\end{table} 

 

\begin{table}[hpt!] 
\begin{center} 
\begin{tabular}{rcccccc} 
\toprule 
$p$ & PCA MV TR  & PCA CP TR & PCA TR & PCA MV ER & PCA CP ER & PCA ER \\ 
\midrule 
$500$  & $3.05$ & $3.49$ & $6.54$ & $3.34$ & $5.16$ & $8.5$ \\ 
$1000$  & $3.07$ & $3.53$ & $6.6$ & $3.31$ & $5.19$ & $8.5$ \\ 
$2000$  & $3.33$ & $3.34$ & $6.68$ & $3.55$ & $4.95$ & $8.5$ \\ 
$3000$  & $3.41$ & $3.29$ & $6.69$ & $3.65$ & $4.85$ & $8.5$ \\ 
\bottomrule 
\end{tabular} 
\end{center} 
\caption{PCA True/Esimtated Return decompositions -- Minimum Variance (MV) and Characteristic Portfolio (CP) -- ($n = 125$, $\mu = 8.5$,  10 simulations).} 
\label{tab:pcaRD} 
\end{table} 

 

\begin{table}[hpt!] 
\begin{center} 
\begin{tabular}{rcccccc} 
\toprule 
$p$ & JSM MV TR  & JSM CP TR & JSM TR & JSM MV ER & JSM CP ER & JSM ER \\ 
\midrule 
$500$  & $3.46$ & $3.39$ & $6.85$ & $4.4$ & $4.1$ & $8.5$ \\ 
$1000$  & $3.87$ & $3.14$ & $7.01$ & $4.76$ & $3.74$ & $8.5$ \\ 
$2000$  & $4.96$ & $2.22$ & $7.18$ & $5.81$ & $2.69$ & $8.5$ \\ 
$3000$  & $5.79$ & $1.49$ & $7.28$ & $6.69$ & $1.81$ & $8.5$ \\ 
\bottomrule 
\end{tabular} 
\end{center} 
\caption{JSM True/Esimtated Return decompositions -- Minimum Variance (MV) and Characteristic Portfolio (CP) -- ($n = 125$, $\mu = 8.5$,  10 simulations).} 
\label{tab:jsmRD} 
\end{table} 

 


\begin{figure}[htp!]
\centering
\includegraphics[width=0.475\textwidth]{img/est_pca_frontier_p_3000.pdf}
\includegraphics[width=0.475\textwidth]{img/act_pca_frontier_p_3000.pdf}
\includegraphics[width=0.475\textwidth]{img/est_mfs_frontier_p_3000.pdf}
\includegraphics[width=0.475\textwidth]{img/act_mfs_frontier_p_3000.pdf}
\caption{The actual (using $\bSig$ and $\alpha$) and estimated 
(using $\hSig$ and $m$) frontiers for the simulation set in 
Table \ref{tab:vol}.}
\end{figure}
\newpage
\section{Recipes}


\begin{enumerate}
\item Starting with $y = B f + \ep$ we first write our model as
\[ y = \alpha  + B x +  u \]
where $x = (f - \Exp(f))$ and $u = \ep - \Exp(\ep)$ so 
that $x,u$ have zero mean and,
\[  \alpha = B \Exp(f) + \Exp(\ep) 
\s[32] \big(\text{and } \Exp(\ep) \perp B \Exp(f) \s . \big)\]

\item From the model above we obtain the $p \times n$
data (for $\nes = (1, \dots, 1) \in \bbR^n$),
\begin{align} \label{uY}
 Y = B \calF^\top + \calE = \alpha \nes_n^\top + 
B \calX^\top + \calU \s .
\end{align}
\item The centered version of the data which we computed by
\begin{align} \label{cY}  
\bar{Y} = Y\s ( \Id_n - \nes\nes^\top /n)
= B \bar{\calX}^\top + \bar{\calU} \s  .
\end{align}
where $\bar{\calX} = \calX ( \Id_n - \nes\nes^\top /n)$
where $\bar{\calU} = \calU ( \Id_n - \nes\nes^\top /n)$. 
\begin{itemize}
\item[--] Letting $\bar{x} = \calX^\top \nes/n$ and 
$\bar{u} = \calU^\top \nes/n$ we see that for each column $x$
and $z$ of $\calX$ and $\calU$ are of the form 
$x - \bar{x}$  and $u - \bar{u}$ 
in the matrices $\bar{\calX}$ and $\bar{\calU}$.
\item[--]  $\Exp(x - \bar{x})$ and 
$\Exp(u - \bar{u})$ are not zero in general! (unless
we have stationary). We assume stationarity 
for $u$ (i.e. $\Exp(u - \bar{u}) = 0$) but not for $x - \bar{x}$.
\item We also have two sample covariance matrices.
\[ S = YY^\top /n \text{ (uncentered)} \s[32]
 \bar{S} = \bar{Y}\bar{Y}^\top /n \text{ (centered)} \]
\end{itemize}
\item The mean vector for the data (that is subtracted during
centering) is, 
\begin{align}\label{by}
 \bar{y} = \frac{Y \nes}{ n \s } 
=  \alpha + B \bar{x} + \bar{u} \end{align}
and our goal is to estimate $\alpha \in \bbR^p$. Our
PCA model will use the constraint, \[ m  = \bar{y} \s . \]
\item To design the JS estimator we follow the standard
playbook. We take
\[ \bar{y}(\rmc) = \rmc \bar{y}  + (1- \rmc) \s g \]
for a shrinkage target $g$. For the grand mean
target, $g = \ip{\bar{y}}{\rme}\s \rme /p$,
%\[ \ip{\bar{y}(\rmc)}{w} = \rmc \ip{\bar{y}}{w}
% + (1 - \rmc) \ip{\bar{y}}{\rme}/p
%= \rmc \ip{m}{w} + (1-\rmc)\ip{m}{\rme} /p \]
\item Our portfolio will hedge out $B \bar{x}$
asymptotically, so we use JS to correct for $\bar{u}$.
Standard calculations that minimize
$|\bar{y}(\rmc)-\alpha + B\bar{x}|^2$ give
the optimal $\rmc$ as
\begin{align}
 \rmc_*
&= \frac{\ip{\bar{y}-g}{\alpha - B \bar{x}-g}}{|\bar{y}-g|^2}
\end{align}
\item Using the LLN behaviour of $\bar{u}$ we have the following
asymptotics,
\begin{align}
\ip{\bar{u}}{\rme}/p, \s[2] \ip{\bar{u}}{\alpha}/p,
\ip{B\bar{x}}{\bar{u}}/p \to 0
\end{align}
\item From the above obtain the asymptotics,
\begin{align}
\rmc_* \sim 1 - \frac{|\bar{u}|^2}{|\bar{y}-g|^2}
\end{align}
The last term (i.e. $|\bar{u}^2|$) is easy to estimate
with eigenvalues (see $\nu^2$ below).
\[  m_{\textsc{js}} = \Big(1 - \frac{\nu^2}{|\bar{y}-g|^2} \Big)
\s  m 
+ \Big(\frac{\nu^2}{|\bar{y}-g|^2} \Big) \s g
\s[32]
\big( g = \ip{\bar{y}}{\rme} \rme / p \big) \]

\item Our MJS estimator below will use $m_\textsc{js}$
in the Markowitz program.

\newcommand{\js}{\ensuremath{\textsc{jsm}}}

\item Write $\bar{S} = HH^\top + G$ be the spectal decomposition 
of the (centered) sample covariance matrix
and let $n_+$ denotes the number of its nonzero eigenvalues,
\[ \seig_{1,p}^2 \ge \seig_{2,p}^2, \dots, \seig^2_{n_+,p} > 0 \s . \]
(If $Y$ is full rank $n_+ = n -1$).
Above $H$ is $p \times q$ matrix of $q$ (spiked) eigenvectors,
\[  H^\top H = \qeig^2 \]
where $\qeig^2$ is the $q \times q$ matrix of spiked
eigenvalues.
\item We write the constraint of the Markowitz program
$\req{mqp}$ in matrix form as
\begin{align} \label{mqp}
 A = \Big(
\begin{array}{cc} m_\js & \rme \end{array}
\Big) \in \bbR^{p \times 2}
\end{align}
\item The average noise in the sample covariance is defined as
\[
 \nu^2 = 
\bigg( \frac{1 + \np}{n_+ - q + \np} \bigg)
\sum_{j > q} \scrs^2_{j,p} \approx
\frac{\sum_{j > q} \scrs^2_{j,p}}{n_+ -q}
\s[32] (n_+ > q) \]
\item We now arrive at our Markowitz-James-Stein estimator.
Let $A^+ = (A^\top A)^{-1} A^\top$.
\begin{align*}
  H^{\text{JSM}} &=  H  C + M  (\Id - C)  \\
 M &=  A A^+ H \\
 N &=   \qeig^2- H^\top AA^+ H \\
C &= (\Id -  N^{-1} \nu^2 )
\end{align*}
\item Alternatively, \begin{align*}
  H^{\text{JSM}} &= H C + M (\Id_q - C) \\
 M &=  A A^+ H \\
 N &= (H- M)^\top (H-M) \\
 C &= \Id - \nu^2 \s N^{-1} 
\end{align*}
\item And alternatively again, for $\scrH = H \calS^{-1}$
(orthonormal columns),\footnote{Brief sketch, We write
$HC + M(\Id - C) = \scrH \qeig  C
+ \scrM \qeig (\Id - C)$ and off we go with
\begin{align*}
\qeig C &= \qeig - \nu^2 \qeig N^{-1} 
= (\Id - \nu^2 \qeig N^{-1} \qeig^{-1}) \qeig 
= (\Id - \nu^2 (\qeig^{-1} N \qeig^{-1})^{-1} \qeig^{-2}) \qeig
= (\Id - \nu^2( N')^{-1} \qeig^{-2}) \qeig  \\
C' &= \Id - (N')^{-1} \nu^2 \qeig^{-2}   \\
N' =  \qeig^{-1} N \qeig^{-1} &=  \Id - \qeig^{-1} H^\top AA^+ H \qeig^{-1}
 =  \Id - \scrH^\top AA^+ \scrH  \\
\qeig (\Id - C) &
= \qeig \nu^2 N^{-1} = \nu^2 (\qeig^{-1}N \qeig^{-1})^{-1}
\qeig^{-1}
= (N')^{-1} \nu^2 \qeig^{-2} \qeig  = (\Id - C') \qeig
\end{align*}
Therefore, $\scrH^{\text{JSM}} = H^\textsc{JSM} \qeig$. }
\begin{align*}
  \scrH^{\text{JSM}} &= \scrH C + \scrM (\Id - C) \\
 \scrM &=  A A^+ \scrH = \scrA \scrA^\top \scrH \\
 N &=  \Id - \scrH^\top AA^+ \scrH 
= \Id - \scrH^\top\scrA\scrA^\top \scrH \\
 C &= I - N^{-1} \nu^2 \qeig^{-2}
\end{align*}
\end{enumerate}

\end{document}

